\section*{अध्याय \ref{ch:g8:speed-of-light} --- प्रकाश की चाल; ब्रह्मांड में दूरियाँ}

\begin{solution}{exo:g8:speed-of-light:1}
$c = \qty{300000}{km/s} = 3 \times 10^{8}\ \unit{m/s}$. पहाड़ियों की
लालटेनों ने सिर्फ़ इंसानी प्रतिवर्त मापे: \omterm{def:g1:light-and-shadows:source}{प्रकाश} ने वे पहाड़ियाँ सेकंड
के दस लाखवें हिस्सों में पार कर लीं, यानी किसी भी हाथ की समय-गिनती से
कहीं नीचे।
\end{solution}

\begin{solution}{exo:g8:speed-of-light:2}
बृहस्पति का \omterm{def:g2:moon-in-the-sky:moon}{चंद्रमा} पाबंद समय-सारणी पर ग्रहण-ग्रस्त होता था --- यानी
आकाश की एक घड़ी। हर ग्रहण की \emph{ख़बर} तब देर से पहुँचती थी जब पृथ्वी
बृहस्पति से दूर खड़ी होती, और पास होने पर जल्दी। उस देरी ने पृथ्वी की
\omterm{def:g6:sun-earth-moon:axisorbit}{कक्षा} की चौड़ाई पार करने में \omterm{def:g1:light-and-shadows:source}{प्रकाश} को लगा समय माप दिया --- और इस तरह
\omterm{def:g1:light-and-shadows:source}{प्रकाश} की \omterm{def:g7:motion-average-speed:formula}{चाल}।
\end{solution}

\begin{solution}{exo:g8:speed-of-light:3}
\omterm{def:g2:moon-in-the-sky:moon}{चंद्रमा}: $384000 \div 300000 \approx \qty{1.3}{s}$. सूर्य:
$150000000 \div 300000 = \qty{500}{s}$ --- आठ मिनट बीस सेकंड।
धूप सूर्य को हमेशा वैसा दिखाती है जैसा वह आठ मिनट पहले था: बनावट से ही
पुरानी ख़बर।
\end{solution}

\begin{solution}{exo:g8:speed-of-light:4}
एक दूरी: वह फ़ासला जो \omterm{def:g1:light-and-shadows:source}{प्रकाश} एक साल में तय करता है --- क़रीब
$9.5 \times 10^{12}$ किलोमीटर।
\end{solution}

\begin{solution}{exo:g8:speed-of-light:5}
चार साल। जवाब अलग-अलग होंगे --- चार साल पहले तुम इसी पुस्तक की कुछ
कक्षाएँ पीछे थे।
\end{solution}

\begin{solution}{exo:g8:speed-of-light:6}
एक ओर के ग्यारह-बारह मिनट के हिसाब से, दिशा सुधारने का कोई भी हुक्म कम
से कम बाईस मिनट पुरानी तसवीर का जवाब देता है --- हत्थे की झटक पहुँचने
से पहले ही रोवर दरार में होगा। इसीलिए योजनाएँ भेजी जाती हैं, हत्थे
नहीं।
\end{solution}

\begin{solution}{exo:g8:speed-of-light:7}
\omterm{def:g2:moon-in-the-sky:moon}{चंद्रमा}: क़रीब \qty{1.3}{s}। सूर्य: आठ मिनट। सबसे पास का \omterm{def:g4:solar-system:star}{तारा}: चार
साल। पड़ोसी मंदाकिनी: पच्चीस लाख साल।
\end{solution}

\begin{solution}{exo:g8:speed-of-light:8}
\omterm{def:g1:light-and-shadows:source}{प्रकाश} पर। \omterm{def:g7:motion-average-speed:formula}{चाल} $c$ धातु की किसी भी छड़ या किसी भी राष्ट्रीय रूलर से
ज़्यादा अटल निकली --- इसलिए अब रूलर \omterm{def:g1:light-and-shadows:source}{प्रकाश} से निकाला जाता है, न कि
\omterm{def:g1:light-and-shadows:source}{प्रकाश} रूलर से मापा जाता है।
\end{solution}

\begin{solution}{exo:g8:speed-of-light:9}
आना-जाना: $300000 \times 2.6 = 780000$ किलोमीटर; और \omterm{def:g2:moon-in-the-sky:moon}{चंद्रमा} उसके आधे
पर खड़ा है --- \qty{390000}{km}। सोनार वाला नियम यहाँ पूरा का पूरा
चलता है: हर \omterm{ex:g4:how-sound-travels:echo}{प्रतिध्वनि} आने-जाने का दाम चुकाती है, इसलिए भरोसा करने से
पहले आधा करो।
\end{solution}

\begin{solution}{exo:g8:speed-of-light:10}
लैंप: $3 \div (3 \times 10^{8}) = 10^{-8}$ सेकंड --- यानी सेकंड के दस
लाखवें हिस्से का सौवाँ भाग: कमरा भी एक अभिलेखागार है, बस इतना उथला कि
ध्यान ही न जाए। फिर सूर्य (आठ मिनट), वह \omterm{def:g4:solar-system:star}{तारा} (एक सदी), और मंदाकिनी
(पच्चीस लाख साल) --- जितना गहरा देखोगे, पन्ना उतना ही पुराना।
\end{solution}

\begin{solution}{exo:g8:speed-of-light:11}
``पिछले मंगलवार किसी तारे के धमाके का \omterm{def:g1:light-and-shadows:source}{प्रकाश} \emph{हम तक पहुँचा} ---
जबकि वह धमाका क़रीब $8000$ साल पहले हुआ था, और उसकी ख़बर तभी से सफ़र
कर रही थी।''
\end{solution}

\begin{solution}{exo:g8:speed-of-light:12}
$3 \times 10^{8}\ \text{km} \div 1000\ \text{s} = 3 \times 10^{5}$
\unit{km/s} --- आधुनिक मान ठीक निशाने पर (और उस ज़माने का असली
अनुमान, \omterm{def:g6:sun-earth-moon:axisorbit}{कक्षा} के मोटे आँकड़ों के साथ, इसके क़रीब चौथाई भर के भीतर बैठा:
सत्रहवीं सदी के लिए हैरतअंगेज़ पहला उत्तर)।
\end{solution}

\begin{solution}{pb:g8:speed-of-light:1}
\textbf{1.} तुम्हारे बोलने के क़रीब \qty{1.3}{s} बाद; और
प्राप्ति-सूचना सबसे जल्दी पूरे आने-जाने के बाद, यानी क़रीब
\qty{2.6}{s} में।
\textbf{2.} एक ओर के लिए $36000 \div 300000 = \qty{0.12}{s}$; और कोई
सवाल-जवाब दो बार ऊपर-नीचे सवारी करता है --- यानी क़रीब आधा सेकंड की
झिझक, ठीक कान की देहरी पर।
\textbf{3.} ``वह लपट हमारे पर्दों के जगमगाने से आठ मिनट बीस सेकंड
पहले फूटी थी --- हम उसे वैसा देखते हैं जैसी वह \emph{थी}।''
\textbf{4.} \omterm{def:g2:moon-in-the-sky:moon}{चंद्रमा} वाली कड़ी: \qty{1.3}{s} बनाम \qty{0.12}{s} ---
यानी मोटे तौर पर दस गुना छूट।
\textbf{5.} $1.8 \times 10^{8} \div (3 \times 10^{5}) = 600$ सेकंड:
यानी एक ओर के दस मिनट।
\textbf{6.} आना-जाना बीस मिनट का है: बचाव वाली तसवीर दस मिनट पुरानी
है और बचाव का हुक्म दस मिनट देर से पहुँचेगा --- दरार उन्नीस मिनट से
जीत जाती है। इसलिए रोवर को अपने प्रतिवर्त साथ लेकर चलने ही पड़ते हैं:
यानी ख़तरा देखकर ख़ुद रुक जाने की व्यवस्था।
\textbf{7.} आख़िरी हिस्सा 21:00 तक पहुँचना चाहिए; वह पृथ्वी से
$21{:}00 - 10$ मिनट $= 20{:}50$ पर चलेगा; और प्रसारण उससे तीन मिनट
पहले शुरू हुआ था: यानी 20:47 पर शुरू करो।
\textbf{8.} देरी दोगुनी होकर एक ओर की क़रीब इक्कीस मिनट हो जाएगी:
ग्रहों तक ``की'' कोई एक दूरी होती ही नहीं --- दोनों दुनियाएँ \omterm{def:g6:sun-earth-moon:axisorbit}{कक्षा}
लगाती हैं, और हर बातचीत एक चलते-फिरते निशाने के हिसाब से तय होती है।
\textbf{9.} $300000 \times 86400 \approx 2.6 \times 10^{10}$
किलोमीटर --- दो अरब साठ करोड़ किलोमीटर; और एक सवाल-जवाब में पूरे
दो दिन लगते हैं।
\textbf{10.} \omterm{def:g1:light-and-shadows:source}{प्रकाश} उस तारे तक क़रीब चार साल में पहुँचता है; और
आज रात के प्रसारण के बाद सबसे जल्दी जवाब क़रीब आठ साल में उतरेगा।
\textbf{11.} ``यह \omterm{def:g1:light-and-shadows:source}{प्रकाश} अपनी मंदाकिनी से पच्चीस लाख साल पहले चला था:
यह तसवीर पत्थर के सबसे शुरुआती औज़ारों वाली पृथ्वी का प्रसारण कर रही
है --- अभिलेखागार का हमारा सबसे गहरा नंगी-आँख वाला पन्ना।''
\textbf{12.} जैसे: ``कुछ भी तत्काल नहीं होता --- हमारे समेत हर संकेत
बहुत हुआ तो $c$ से चलता है। \omterm{def:g2:moon-in-the-sky:moon}{चंद्रमा} से आगे हत्थे देरी से हार जाते हैं:
दूर की मशीनों को ख़ुद सोचना पड़ता है। और हर दूरबीन अभिलेखागार का पाठक
है --- वस्तु जितनी दूर, ख़बर उतनी पुरानी, और कोई ताज़ा संस्करण छापा ही
नहीं जा सकता।''
\end{solution}
